Name der Einrichtung;Strasse mit Hausnummer;Postleitzahl;Stadt;E-Mail-Adresse
SafeZone.ch Online-Beratung zu Suchtfragen www.safezone.ch Portal zu Suchtfragen mit Beratungen, Selbsttests, Suchthilfeangeboten und Substanzwarnungen. Kostenlos und anonym Details für «SafeZone.ch» anzeigen;;;;
Kantonsspital Aarau, Pneumologie und Schlafmedizin Rauchstopp-Sprechstunde Tellstrasse 25, 5001 Aarau www.ksa.ch rauchstoppberatung@ksa.ch 062 838 44 78 Rauchen führt zu einer körperlichen und psychischen Abhängigkeit. Für diejenigen, welche mit dem Rauchen aufhören wollen, stellen diese Abhängigkeiten eine zusätzliche Hürde dar, die überwunden werden muss. Wir bieten eine Beratung und Unterstützung beim Rauchstopp an, bei welcher wir nicht nur Medikamente gegen die körperliche Abhängigkeit und die Lust auf die Zigarette anbieten, sondern Rauchern auch Strategien aufzeigen, wie sie die Gewohnheiten des Rauchens durchbrechen und durch alternative Handlungen ersetzten können. Mittels Messung des ausgeatmeten CO sowie mittels Spirometrie lässt sich rasch eine Verbesserung der Gesundheit nach dem Rauchstopp dokumentieren. Details für «Kantonsspital Aarau, Pneumologie und Schlafmedizin» anzeigen;Kantonsspital Aarau, Pneumologie und Schlafmedizin Rauchstopp-Sprechstunde Tellstrasse 25;5001;Aarau www.ksa.ch rauchstoppberatung@ksa.ch 062 838 44 78 Rauchen führt zu einer körperlichen und psychischen Abhängigkeit. Für diejenigen;rauchstoppberatung@ksa.ch
"Lungenliga Aargau Rauchstopp Hintere Bahnhofstrasse 8, 5001 Aarau www.lungenliga.ch rauchstopp@llag.ch 062 832 40 14 Angebote: - Infoabend Rauchstopp - Persönliches Rauchstopptraining für Erwachsene und Jugendliche - Rauchstoppkurs «ich schaff&apos;s!» - WhatsApp-Chat Nikotinstopp Details für «Lungenliga Aargau» anzeigen";Lungenliga Aargau Rauchstopp Hintere Bahnhofstrasse 8;5001;"Aarau www.lungenliga.ch rauchstopp@llag.ch 062 832 40 14 Angebote: - Infoabend Rauchstopp - Persönliches Rauchstopptraining für Erwachsene und Jugendliche - Rauchstoppkurs «ich schaff&apos;s!» - WhatsApp-Chat Nikotinstopp Details für «Lungenliga Aargau» anzeigen";rauchstopp@llag.ch
Suchtberatung ags, Aarau Metzgergasse 2, 5000 Aarau www.suchtberatung-ags.ch aarau@suchtberatung-ags.ch 062 837 60 40 Information, Beratung, Kurzberatung per E-Mail, psychosoziale Begleitung, Krisenintervention, Vermittlung von Behandlungsangeboten (Entzug, stationäre Therapie, Substitutionsbehandlung, Medikamente usw.). Abgabe von sterilem Spritzenmaterial und Kondomen (Montag bis Donnerstag 13.30 - 16.30 Uhr, Freitag 13.30 - 16 Uhr). Glücksspiel- und Gamesuchtberatung. Erstabklärung und Vermittlung an zuständige Stellen bei Essstörungen und für suchtmedizinische Behandlungen. Details für «Suchtberatung ags, Aarau» anzeigen;Suchtberatung ags, Aarau Metzgergasse 2;5000;Aarau www.suchtberatung-ags.ch aarau@suchtberatung-ags.ch 062 837 60 40 Information;aarau@suchtberatung-ags.ch
Selbsthilfegruppen Blaubrügg Blaues Kreuz Bern-Solothurn-Freiburg 3715 Adelboden www.besofr.blaueskreuz.ch selbsthilfe@blaueskreuzbern.ch 033 222 01 77 Hilfe zur Selbsthilfe. Die Selbsthilfegruppen Blaubrügg bieten Menschen mit Sucht- und Alkoholproblemen und ihren Angehörigen eine Möglichkeit zum Erfahrungsaustausch. Sie stossen auf Verständnis, unterstützen sich gegenseitig und hören einander zu. Niemand wird verurteilt, es werden andere Lösungswege aufgezeigt und Therapieerfolge mit einbezogen. Im Kontakt mit Gleichbetroffenen wird die Isolation durchbrochen. Die Treffen gestaltet jede Gruppe selbst. Studien zeigen: Die Teilnahme in einer Selbsthilfegruppe kann ein wichtiger Schutzfaktor gegen Rückfälle sein. Es bestehen Gruppen in Bern, Biel, Thun und Adelboden. Gerne unterstützen wir neue Selbsthilfegruppen im Aufbau. Details für «Selbsthilfegruppen Blaubrügg» anzeigen;;;;selbsthilfe@blaueskreuzbern.ch
Kanton Zürich, Amt für Jugend und Berufsberatung CONTACT Jugendberatung Bezirk Affoltern Im Winkel 2, 8910 Affoltern am Albis www.contact-jugendberatung.ch info@contact-jugendberatung.ch 043 259 93 55 Die Jugendberatung CONTACT unterstützt Jugendliche und junge Erwachsene im Bezirk Affoltern auf dem Weg in ein selbständiges Leben. Auch Eltern, Lehrpersonen oder Arbeitgeber können sich an uns wenden. Wir beraten bei Fragen zu körperlichen und psychischen Veränderungen in der Pubertät, zu Sexualität und Liebesbeziehungen oder zum Umgang mit Genuss- und Suchtmitteln. Details für «Kanton Zürich, Amt für Jugend und Berufsberatung» anzeigen;Kanton Zürich, Amt für Jugend und Berufsberatung CONTACT Jugendberatung Bezirk Affoltern Im Winkel 2;8910;Affoltern am Albis www.contact-jugendberatung.ch info@contact-jugendberatung.ch 043 259 93 55 Die Jugendberatung CONTACT unterstützt Jugendliche und junge Erwachsene im Bezirk Affoltern auf dem Weg in ein selbständiges Leben. Auch Eltern;info@contact-jugendberatung.ch
"Sozialdienst des Bezirks Affoltern, Suchtberatung Zweckverband Sozialdienst des Bezirks Affoltern Obfelderstrasse 41b, 8910 Affoltern am Albis www.sozialdienst-affoltern.ch info@sdaffoltern.ch 044 762 45 45 Der Sozialdienst des Bezirks Affoltern ist eine Institution der im Zweckverband ""Sozialdienst des Bezirks Affoltern"" zusammengeschlossenen Gemeinden des Bezirks Affoltern. Wir geben Auskunft, Rat und Hilfe bei Suchtgefährdung, Abhängigkeit und Missbrauch von legalen und illegalen Suchtmitteln sowie bei abhängigem Verhalten. - Einzel-, Paar- und Familiengespräche - Informationen über Entzug, stationäre Therapieeinrichtungen sowie Nachfolgeprogramme/-behandlungen - Abklärung für Eintritt in ambulante, halbstationäre oder stationäre Entzugs- oder Langzeittherapien - Begleitung und Beratung von Personen nach einem stationären Aufenthalt in ihrer Stabilisierungs- oder Reintegrationsphase Details für «Sozialdienst des Bezirks Affoltern, Suchtberatung» anzeigen";Sozialdienst des Bezirks Affoltern, Suchtberatung Zweckverband Sozialdienst des Bezirks Affoltern Obfelderstrasse 41b;8910;"Affoltern am Albis www.sozialdienst-affoltern.ch info@sdaffoltern.ch 044 762 45 45 Der Sozialdienst des Bezirks Affoltern ist eine Institution der im Zweckverband ""Sozialdienst des Bezirks Affoltern"" zusammengeschlossenen Gemeinden des Bezirks Affoltern. Wir geben Auskunft";info@sdaffoltern.ch
Lungenliga Zentralschweiz, Beratungsstelle Altdorf Rauchstopp Spitalstrasse 1a, 6460 Altdorf www.lungenliga.ch info@lungenliga-zentralschweiz.ch 041 429 31 10 Die Lungenliga Zentralschweiz deckt folgende Kantone ab: LU, NW, OW, SZ, UR, ZG Angebote: - Einzelberatung telefonisch oder persönlich - Kostenlose Rauchstopp-Beratung für Jugendliche, einzeln oder in Kleingruppen - Kurzintervention telefonisch - Rauchstopptrainings in Betrieben Details für «Lungenliga Zentralschweiz, Beratungsstelle Altdorf» anzeigen;Lungenliga Zentralschweiz, Beratungsstelle Altdorf Rauchstopp Spitalstrasse 1a;6460;Altdorf www.lungenliga.ch info@lungenliga-zentralschweiz.ch 041 429 31 10 Die Lungenliga Zentralschweiz deckt folgende Kantone ab: LU;info@lungenliga-zentralschweiz.ch
Stiftung Papilio Jugend- und Elternberatung, Suchtberatung Gotthardstrasse 14, 6460 Altdorf www.stiftung-papilio.ch info@stiftung-papilio.ch 041 874 13 29 Wir beraten Betroffene, deren Angehörige und Bezugspersonen. Das Tätigkeitsgebiet unserer Beratung umfasst sowohl legale (Alkohol, Medikamente, Nikotin) als auch illegale (Cannabis, Kokain, Heroin usw.) und stoffungebundene Suchtformen (Spielsucht, Computersucht etc.). Details für «Stiftung Papilio» anzeigen;Stiftung Papilio Jugend- und Elternberatung, Suchtberatung Gotthardstrasse 14;6460;Altdorf www.stiftung-papilio.ch info@stiftung-papilio.ch 041 874 13 29 Wir beraten Betroffene;info@stiftung-papilio.ch
Zentrum Selbsthilfe Uri Gotthardstrasse 14, 6460 Altdorf UR www.selbsthilfe-uri.ch info@selbsthilfe-uri.ch 0418741394 Anlauf-, Informations- und Koordinationsstelle rund um das Thema Selbsthilfe und Selbsthilfegruppen im Kanton Uri. Details für «Zentrum Selbsthilfe Uri» anzeigen;Zentrum Selbsthilfe Uri Gotthardstrasse 14;6460;Altdorf UR www.selbsthilfe-uri.ch info@selbsthilfe-uri.ch 0418741394 Anlauf-;info@selbsthilfe-uri.ch
